\documentclass[12pt, a4paper]{achemso}
% Authors
\author{Elvis A. F. Martis}
\affiliation[Nantes University]
{US2B, Nantes Université, CNRS, UMR 6286, F-44000 Nantes, France}
\author{Johann Hendrickx}
\affiliation[Nantes University]
{US2B, Nantes Université, CNRS, UMR 6286, F-44000 Nantes, France}
\author{Vincent Sauzeau}
\affiliation[CHU Nantes] {UMR\_S 1087 {/} UMR\_C 6291 l'Unité de Recherche de l'Institut du Thorax, F-44000 Nantes, France}
\author{Stéphane Téletchéa}
\email{stephane.teletchea@univ-nantes.fr}
\affiliation[Nantes University]
{US2B, Nantes Université, CNRS, UMR 6286, F-44000 Nantes, France}

% Title
\title
{Detection of Alternate Druggable Pockets on RAC1}


\begin{document}
 % Abstract
 \begin{abstract}
  RAC1 (Ras-related C3 botulinum toxin substrate 1) is a molecular switch catalysing by hydrolysing GTP to GDP. Owing to the slow nature of GTP hydrolysis and even slower dissociation of GDP, RAC1 is always assisted by Guanine Exchange Factors (GEFs) and GTPase-activating proteins (GAPs). As a result, GEFs and GAPs control the switching mechanism in RAC1. This regulatory control mechanism by several GEFs and GAPS provides a means to modulate the RAC1 activity in several diseases in which it has been implicated. The work described here attempts to understand various dynamical aspect of RAC1 in presence of GTP, GDP and guanine exchange factors (GEFs). The primary objective is to enhance the scope of therapeutic interventions in RAC1-associated disease pathways by locating alternative binding site(s) in the RAC1 surface not observed in X-ray derived tertiary structure.  We are aware that targeting the nucleotide binding pocket is challenging due to the picomolar binding affinity and high levels of GTP to compete with the drug molecules. we embarked on finding alternative pockets on RAC1 that can modulate its activity. We observe that the active conformation of RAC1, GTP- or GEF-bound, opens a cryptic pocket that lies below the switch I region. This pocket could be targetted by small molecules to control the flexibility of switch I, which is crucial for nucleotide exchange, and therefore, for GEF binding. 
\end{abstract}

% Main Text starts from here

% Introduction and Motivation of the work
\section{Introduction}
    Small GTPases, a ubiquitous class of proteins, play pivotal roles in various cellular processes. Their functions include signal transduction, vesicle trafficking, cytoskeletal dynamics, and cell division. Ras GTPases regulate cell proliferation, differentiation, and survival. While Rho GTPases govern cytoskeletal dynamics and cell migration.  Their activity is regulated through a molecular switch mechanism by guanine nucleotide exchange factors (GEFs), GTPase-activating proteins (GAPs), and guanine nucleotide dissociation inhibitors (GDIs)\cite{cherfils2013regulation}. These molecular switches cycle between OFF (GDP-bound inactive) and ON (GTP-bound active) states. This switching mechanism controls various processes like cell growth and differentiation to vesicular and nuclear transport\cite{cherfils2013regulation, loirand2013small}. 
    Small GTPases are proteins that possess conserved domains across several families; made up of six $\beta$-sheets surrounded by 5 $\alpha$-helices. All domains are conserved such that they carry a guanine nucleoside binding site divided to recognize the guanine base and phosphate groups and the magnesium ion\cite{toma2019structural, cherfils2013regulation}. The switch regions regulate the traffic to the guanine nucleoside binding site. It is also well established that GTPases alone cannot hydrolyse GTP effectively. Therefore,  GEFs and GAPs play a significant role in assisting the switching mechanism of GTPases. They control the dynamics of the switch regions that facilitate either hydrolysis of GTP or dissociation of GDP. GDP:Mg$^{2+}$ dissociates from GTPases very slowly, and therefore, GEFs assist in exchanging GDP for GTP One of the accepted mechanisms of GEFs in the dissociation of GDP is by electrostatic repulsion of magnesium ion that forms strong interaction with GTPases and phosphate groups of GDP. The loss of magnesium ion weakens the affinity of GDP which eventually dissociates. Alternatively, some GEFs can rearrange the switch I or II conformations such that GDP's affinity is reduced facilitating its dissociation. However, GEFs can also combine the two mechanisms to facilitate GDP dissociation \cite{cherfils1999gefs, cherfils2014gefs}. \\ 
    Furthermore, GAPs stabilise the transition state for GTP:Mg$^{2+}$:GTPases and facilitate the hydrolysis of $\gamma$-phosphate group from GTP. taken together, the switching mechanism can be seen as GEFs assist in the GTPase activation while GAPs assist in GTP hydrolysis. Following GTP hydrolysis, the inactive GTPases prevent further signal transduction.\cite{bos2007gefs}. 
    In some classes of GTPases, the active and inactive state cycling is conjugated with GTPases being shuttled between the cell membrane and cytosol. Such GTPases are found with post-translational modification to carry long-chain fatty acids at their C-terminal to assist bilayer insertion. GDIs\cite{dermardirossian2005gdis, maruta1998gtpase} sequester the lipid component to prevent membrane insertion to GTPases in the cytosol. This coordinated action of GTPases, GEFs, GAPs, and GDIs is essential for various cellular processes. Dysregulation of any of these components contributes to pathologies like cancer, neurodegeneration, and inflammatory disorders. For instance, mutations in Ras and its regulators are implicated in oncogenesis, highlighting the therapeutic potential of targeting GTPase signalling pathways\cite{cherfils2013regulation}. Such mutations are known to disrupt the equilibrium between the active and inactive states of GTPases. It has also been observed that some mutations enhance the affinity of GEFs or GAPs that traps the small GTPases either in an activated state or inactivated state. Small molecule inhibitors targeting GTPase activity or disruptors of GEF/GAP interactions\cite{gray2020targeting} are under development for cancer therapy. However, there are several challenges in targeting GTPases and/or GTPases-GAP/GEF complexes. Firstly, targeting the guanine binding site is challenging owing to the large intracellular concentration of GTP and high affinity\cite{traut1994physiological}. In such cases, finding an allosteric site can avoid such competition with endogenous substrates. GTPase hyperactivation or hyper-inactivation can be corrected by disrupting the protein-protein interaction (PPI) between GTPases and GEFs/GAPs. However, it is known that targeting PPI interfaces is extremely challenging due to shallow cavities. 
   

% Methods section
\section{Methods}

\begin{table}
  \caption{Detailed system studies}
  \label{tbl: Detailed system studies}
  \begin{tabular}{llccc}
    \hline
   \textbf{Initial PDB} & \textbf{Protein} & \textbf{Modulator/LIGAND} & \textbf{Size (nb atoms)} \\
    \hline
    1HE1\cite{gappdb} & RAC1 (INACTIVE) & GDP:$Mg^{2+}$ & 26102 \\
    3TH5\cite{activeRAC1}   & RAC1 (ACTIVE) &  GTP:$Mg^{2+}$ & 24132  \\
    4YON\cite{prex1pdb} & RAC1:PREX1    & -- & 88951 \\
    1FOE\cite{tiam1pdb} & RAC1:TIAM1     & -- & 64367 \\
    2VRW\cite{vav1pdb} & RAC1:VAV1      & -- & 61807 \\
    2YIN\cite{dock2pdb} & RAC1:DOCK2      & -- & XXXX \\
    \hline
  \end{tabular}
  
  \textsuperscript{\emph{1}}See the supporting information
\end{table}


\begin{itemize}
    \item List of PDBS used. All systems modelled are described in Table \ref{tbl: Detailed system studies}. 
\end{itemize}

\subsection{Preparation of systems for molecular dynamics simulation:} 
    The coordinates for the protein-ligand complexes were retrieved from the Protein Data Bank (PDB); the PDB entries used in this study are listed in \textbf{Table \ref{tbl: Detailed system studies}}. The X-ray structures were imported and processed in the UCSF Chimera v1.17 \cite{pettersen2004ucsf}. Any missing loops were built using modeller \cite{fiser2003modeller} with  UCSF Chimera v1.17 as the GUI \cite{yang2012ucsf, huang2014enhancing}. 


\subsection{Generating AMBER-compatible parameters:} 
 All systems were prepared for MD simulations with the Leap module in AMBERtools23. Protein atoms were atom-typed according to the AMBER14 (ff14SB) force field\citation{ff14sb}.  Amber-compatible parameters for all nucleotides were obtained from the Manchester AMBER database  For the VAV1 zinc finger region the parameters of the Zinc atom, coordinating cysteines and histidine residues, the force field parameters were obtained from the ZAFF\cite{zaff} force field library.

\subsection{Molecular dynamics simulations:} 
 All complexes were solvated with TIP3P\cite{tip3p} waters enclosed in a box with dimensions extending 10\r{A} beyond any protein atom,  and an appropriate number of neutralizing ions were added to maintain the system's neutrality.  The complexes were relaxed using three cycles of minimization of 50000 steps each.  The first 5000 steps of minimization adopted the steepest descent algorithm, and the remaining 45000 steps used the conjugate gradient algorithm.  The cut-off for evaluating non-bonded interactions was placed at 8.0\r{A} for the minimization phase.  The first cycle of minimization was performed with  25 kcal/(mol{\r{A}$^2$} restraining force on the solute atoms, this force constant was reduced to 5 kcal/(mol\r{A}$^2$)  in the next cycle, and the final cycle was performed without any restraining forces on either the solute or the solvent atoms.  The system was heated in 4 stages (see \textbf{Table \ref{tbl: Heating using NVT Ensemble}}) from 0 K to 300 K for 10 ns with  50 kcal/(mol{\r{A}$^2$} restraining force on the solute atoms using the Langevin dynamics, with a collisional frequency of  2 $ps^-1$.  The next step employed the NPT ensemble for density equilibration for another 2 cycles of 10 ns each, with the restraining force on the solute atoms gradually decreasing from 5 to 0 kcal/(mol{\r{A}$^2$}.  In the equilibration phase, the cut-off for evaluating non-bonded interactions was placed at 8 {\r{A}}.  In the production phase of the MD simulations, all systems were simulated for 100 ns in triplicate, amounting to a 300 ns trajectory for each system.  The long-range electrostatic interactions were calculated using the Particle Mesh Ewald (PME) summation method\cite{pme}. The cut-off for evaluating non-bonded interactions was placed at 9.0 {\r{A}}  for the production phase. All the computations were done using the CUDA-enabled PMEMD code\cite{pmemd01, pmemd02, pmemd03} in AMBER22\cite{amber22, amber}.  
All simulations were performed with the SHAKE algorithm\cite{shake} to constrain bonds to hydrogen atoms, with the tolerance set to 0.00001 {\AA}.  A 2 fs integration time step was used to solve Newton\textquotesingle s equations of motion.  For analysis, the conformations were saved every 100 ps, amounting to 1000 snapshots per trajectory (3000 per system). 

\begin{table}
  \caption{Details of stages of Heating using NVT Ensemble}
  \label{tbl: Heating using NVT Ensemble}
  \begin{tabular}{lcccc}
    \hline
   \textbf{Stage}  & \textbf{Temperature range (K)} & \textbf{Time (ns)} \\
    \hline
    Stage 1 & 0 to 100 & 0 to 2.5 \\
    Stage 2 & 100 to 200 & 2.5 to 5 \\
    Stage 3 & 200 to 300 & 5 to 7.5 \\
    Stage 4 & 300 to 300 & 7.5 to 10 \\ 
    \hline
  \end{tabular}
\end{table}

\subsection{Measuring the structural dynamics of RAC1:} 
    Analysis of the MD trajectory was performed with the cpptraj\cite{cpptraj} module in AMBERTools23.  Before the analysis, all solvent atoms and neutralizing ions were stripped off, and the protein conformations were aligned to their respective X-ray conformation.  The mass-weighted C$\alpha$ root-mean-square deviation (RMSD) for the protein was computed using the X-ray conformation as the reference. The root-mean-square fluctuations (RMSF) per residue for the backbone and side-chain atoms were computed to understand the domain-specific fluctuations. 
    Furthermore, to understand the specific dynamical changes in RAC1 when bound to different binding partners, we extracted the essential dynamical properties using Principal Component Analysis (PCA). The Variance-Covariance matrix for C$\alpha$ atoms was extracted in  cpptraj\cite{cpptraj} module in AMBERTools23. We extracted and analysed 10 principal components, however, the first two principal components explained $<$ 50\% variance in the data. To visualise the essential dynamics, the vectors were projected on the C$\alpha$ atoms and visualised in VMD 1.9.3 \cite{humphrey1996vmd}.

\subsection{Identification and measurement of alternate pockets on RAC1:} 
    POVME2 \cite{durrant2011povme, durrant2014povme} was employed for identifying various pockets on the RAC1. To be able to incorporate various experimental conformations of RAC1, we employed RAC1 with different binding partners (see \textbf{Table \ref{tbl: Detailed system studies}}).  To identify the pockets, the pocketID module on POVME2 was used on aforementioned RAC1 structures. The parameters used for pocket identification are shown in \textbf{Table~\ref{tbl:pocketID}}. Following the pocket identification, we monitored the pocket volume changes of these pockets throughout MD trajectories. It has been observed that the pocket might be a good candidate for druggability if its volume range\cite{perot2010druggable, nayal2006nature, an2005pocketome} was found to lie between 610\r{A}$^3$ to 930\r{A}$^3$. 
    
    \begin{table}
  \caption{Parameters for Pocket Identification using POVME-2}
  \label{tbl:pocketID}
  \begin{tabular}{lc}
    \hline
   \textbf{Parameter}  & \textbf{Value} \\
    \hline
    pocket detection resolution	 & 4.0 \r{A}  \\
    pocket measuring resolution & 1.0 \r{A} \\
    clashing cutoff & 3.5 \r{A}  \\
    number of neighbors & 4  \\ 
    number of spheres & 5  \\
    sphere padding & 5.0 \r{A}  \\ 
    \hline
  \end{tabular}
\end{table}

\begin{itemize}
    \item Analysis of other MD observations to support pocket volume changes
        \begin{itemize}
            \item RMSDs for switch regions 1 and 2
            \item (perhaps) PCA analysis to explain the switch region dynamics and the pocket volumes
        \end{itemize}
\end{itemize}


\subsection{Quantifying the GEF binding hotspots on RAC1 surface:} 

The free energy of binding was computed using the equation (Eq-\ref{eq:free-energy}) on the conformations saved from the MD simulations using mm-pbsa python script\cite{mmpbsa} available in AMBERTools23; 

\begin{equation} \label{eq:free-energy}
 \   {\Delta}G_{binding} = {\Delta}E_{MM} + {\Delta}H_{GBSA} - T{\Delta}S   
\end{equation}

Where  \({\Delta}E_{MM}\) is the interaction energy between the binding partners computed in the gas phase using a molecular mechanics method.  The solvation-free energy or the enthalpy of hydration \({\Delta}H_{GBSA}\)  was computed using the Generalized-Born Surface Area (GBSA) method, and the entropic component of binding \(T{\Delta}S\) was not computed due to huge computation cost. The internal dielectric constant for GBSA calculations was set to 1 and the solvent dielectric constant was fixed as 80. For computing free energy of binding using the MM-GBSA method, the OBC’s GB model\cite{igb5} $ (igb = 5) $ was used. The mbondi3 radii set was employed to define the PB radii for all solute atoms during the preparation of the topology and parameter files in the leap module in AMBERTools23. The ionic strength for  GBSA calculations was maintained at 150 mM. All energy components, in an implicit solvent, were calculated using a large cut-off of 999 Å. The surface area for the non-polar solvation term was computed using the linear combination of pairwise overlaps (LCPO) method\cite{lcpo}. To identify the critical amino acid residues that are critical for the protein-protein interactions between RAC1 and GEFs, we computed the pairwise decomposition of the binding affinity value. 


% Results and Discussion
\section{Results and discussion}

\subsection{Nucleotide and GEF binding significantly affects the switch region dynamics}
 Comments on RMSD
    a) Region wise w/ and w/o nucleotides
    b) Region wise w/ and w/o GEFs
    c) Overall comment on a) and b)
    d) Show with RMSF plots the flexibility of each region 
    e) Support with PCA plots for more detailed discussion (later this will support different pocket dynamics)

\subsection{identifying alternate binding pockets on RAC1 surface} 
    I) Discuss elaborately the RAC1 regions that significantly contribute to GEF binding vis-à-vis X-ray observations on these interactions. 
        a) Link and compare to the dynamical observation of these regions/residues w/o GEFs.
        

\bibliography{References} 

\end{document}